% generator_I35.tex -- Prop I.35 (Theor.): parallelograms on same base between same parallels are equal in area.
\documentclass[12pt]{article}\usepackage{geometry}\geometry{a5paper,margin=1.5cm}
\usepackage[lining]{ebgaramond}\usepackage{amsmath}\usepackage{amssymb}\usepackage{byrne}
\def\mpPre{u := 1cm; textLabels := false;}
\begin{document}
\defineNewPicture[1/5]{%
  pair A,B,C,D,E,F,G,d[];%
  d1:=(7/4u,0);d2:=(u,-3u);d3:=(-2u,-3u);%
  A:=(0,0);B:=A shifted d1;C:=A shifted d2;D:=C shifted d1;%
  E:=C shifted -d3;F:=D shifted -d3;%
  G:=(B--D) intersectionpoint (C--E);%
  draw byPolygon(A,B,G,C)(byyellow);%
  draw byPolygon(E,F,D,G)(byyellow);%
  draw byPolygon(B,E,G)(byyellow);%
  draw byPolygon(C,D,G)(byblue);%
  byAngleDefine(E,A,C,byred,SOLID_SECTOR);%
  byAngleDefine(F,B,D,byblue,SOLID_SECTOR);%
  byAngleDefine(A,E,C,byblack,SOLID_SECTOR);%
  byAngleDefine(B,F,D,white,SOLID_SECTOR);%
  draw byNamedAngleResized();%
  draw byNamedAngleDummySides(BFD);%
  byLineDefine(A,C,byblue,SOLID_LINE,REGULAR_WIDTH);%
  byLineDefine(B,D,byred,SOLID_LINE,REGULAR_WIDTH);%
  byLineDefine(C,D,byblack,SOLID_LINE,REGULAR_WIDTH);%
  draw byNamedLineSeq(0)(AC,CD,BD);%
  draw byLabelsOnPolygon(C,A,B,E,F,D)(ALL_LABELS,0);%
}
\noindent\begin{minipage}[t]{0.45\linewidth}\centering\vspace{0pt}%
\drawCurrentPicture\end{minipage}\hfill
\begin{minipage}[t]{0.52\linewidth}\vspace{0pt}%
\textbf{Book~I.\enspace Prop.\ XXXV. Theor.}
\medskip\noindent\itshape
Parallelograms on the same base, and between the same
parallels, are (in area) equal.
\bigskip\noindent\upshape\begin{center}
On account of the parallels,\\
$\left.\begin{aligned}
\drawAngle{A}&=\drawAngle{B};\\
\drawAngle{E}&=\drawAngle{F};\\
\mbox{and }\drawUnitLine{AC}&=\drawUnitLine{BD}
\end{aligned}\right\}$~(pr.~I.29,~I.34).\\[6pt]
But $\drawPolygon{ABGC,BEG}=\drawPolygon{EFDG,BEG}$~(pr.~I.8)\\[4pt]
$\therefore \drawPolygon{ABGC,EFDG,BEG,CDG}-\drawPolygon{EFDG,BEG}$\\
$\quad=\drawPolygon{ABGC,CDG}$,\\
and $\drawPolygon{ABGC,EFDG,BEG,CDG}-\drawPolygon{ABGC,BEG}$\\
$\quad=\drawPolygon{EFDG,CDG}$;\\[4pt]
$\therefore \drawPolygon{ABGC,CDG}=\drawPolygon{EFDG,CDG}$.
\end{center}
\begin{flushright}Q.\ E.\ D.\end{flushright}
\end{minipage}
\end{document}
